\documentclass[11pt]{article}
\usepackage[a4paper,margin=27mm]{geometry}
\usepackage[T1]{fontenc}
\usepackage{lmodern,microtype,amsmath,amssymb,amsthm,enumitem,booktabs}
\usepackage[hidelinks]{hyperref}
\newtheorem{proposition}{Proposition}
\newtheorem{lemma}[proposition]{Lemma}
\theoremstyle{remark}\newtheorem{remark}[proposition]{Remark}
\newcommand{\Dscr}{\mathcal D}
\newcommand{\C}{\mathcal C}
\newcommand{\M}{\mathcal M}
\title{Vacuous Targets and Transfinite Construction}
\author{}\date{12 October 2026}
\begin{document}\maketitle

\noindent
The new complete-ranking impossibility theorem settles full coverage
under strict guidance improvement, independence, and closed competitive
selection. It does not settle which proper credal targets these rankings
can cover. This note separates that issue from the treatment of the
vacuous target, and records what a transfinite construction would still
have to prove. The two elementary constructions and the union lemma
below are author-checked, not independently refereed.

\section{What the new theorem establishes}

\emph{Three Complete Rankings Cannot Cover} proves that complete
rankings satisfying (G), (IB), and (S) dominate the vacuous model $V$.
\emph{Complete Rankings for Any Number of Worlds} extends this to every
finite $n\ge3$. The counting lemma is the key new step: it uses only
transitivity, weak (G), and (IB), and bounds the number of disjoint heavy
sets using a finite connecting chain. Closed competitive selection
provides the chains and the light perturbations; completeness is used
in obtaining them and in converting failed comparisons into strict
reverse comparisons.

The conclusion concerns a class of rankings as follows. Axiom (A)
requires, for each coherent model, \emph{some} ranking system under which
it is undominated. It does not require one system under which all
coherent models are undominated. Nevertheless, if every permitted
system has complete rankings and satisfies strict (G), (IB), and (S),
none can protect $V$. Therefore the class cannot satisfy full (A).

This is a genuine obstruction to full (A), but is not an impossibility
theorem for imprecision in general. For $n\ge4$, the many-world note
also gives a conditional result for targets whose credal sets contain
three vertices; its universal domination claim for all such targets
remains open. For $n=3$ the only closed convex credal set containing all
three vertices is $\Delta$ itself.

\section{Strict guidance improvement can protect the vacuous model}

Use the papers' binary-problem space $H$, with finite Gaussian measure
$\mu$, and identify models modulo null sets. A problem is \emph{mixed}
if the worlds disagree about which option is better, and
\emph{unanimous} otherwise. The model $V=\Dscr_\Delta$ gives silence on
mixed problems and the unanimously good recommendation on unanimous
ones. Advice at world $i$ is good, silent, bad, or contradictory.

\begin{proposition}\label{prop:strict}
For every finite $n\ge2$, there are complete real-valued additive rankings
on all measurable models satisfying strict (G) and (IB) under which $V$
is undominated, indeed the unique minimizer of the sum of the losses.
\end{proposition}
\begin{proof}
Choose $0<c<1/n$. At every problem and every world assign penalties
\[
 \begin{array}{c|cccc}
 \text{advice}&\text{good}&\text{silence}&\text{bad}&\text{contradictory}\\
 \hline
 \text{penalty}&0&c&1&2
 \end{array}
\]
and let $L_i(X)$ be the integral of the penalty at world $i$.
The penalties are strictly increasing in guidance rank, so the rankings
satisfy strict (G). Additivity makes replacement of common advice
irrelevant to a comparison, so (IB) holds.

On a mixed problem, silence has total penalty $nc<1$. Either definite
recommendation is bad at some world and consequently has total penalty
at least $1$. Contradiction has total penalty $2n$. Silence therefore
uniquely minimizes the total. On a unanimous problem the unanimously
good recommendation uniquely minimizes the total, at zero. Integrating
shows that $V$ uniquely minimizes $\sum_iL_i(X)$ modulo null sets.
A dominator would lower that sum, which is impossible.
\end{proof}

Thus giving good advice can have strictly positive value over silence
even when $V$ is undominated. The decision-theoretic interpretation is
ordinary caution: an error costs enough that abstention can be optimal
under uncertainty. For three worlds, for example, take $c=1/4$.
Silence costs $3/4$ across the worlds, whereas either recommendation on
a mixed problem costs at least $1$.

For $n\ge3$ these rankings cannot satisfy (S), by the new theorem.
The incompatibility concerns the full conjunction of requirements, not
strict (G) by itself. A related stronger example already appears in
\emph{Beyond Additive Guidance Rankings}, in its section on completeness
without (S): complete independent rankings can leave every coherent
model undominated when (S) is omitted.

\section{A precise way to allow neutral abstention}

One may nevertheless want a ranking tailored to $V$ to give no reward
for choosing a good option on a problem about which the worlds disagree.
The following version keeps a strict incentive to follow unanimous
advice. That qualification avoids tying $V$ with the incoherent model
that gives no advice even on unanimous problems.

\begin{proposition}\label{prop:neutral}
There are complete additive rankings on all measurable models satisfying
(IB), weak (G), (E), (LA), (S), and closure under equal undominated
profiles (C), whose only undominated model is $V$, modulo null sets.
The sole failure of the strict guidance requirement is that good advice
and silence tie on mixed problems.
\end{proposition}
\begin{proof}
Integrate the following penalties at each world:
\[
 \begin{array}{c|cccc}
 &\text{good}&\text{silence}&\text{bad}&\text{contradictory}\\
 \hline
 \text{mixed problem}&0&0&1&2\\
 \text{unanimous problem}&0&1&2&3
 \end{array}
\]
All losses are finite and nonnegative; $L_i(V)=0$ at every world.
If a model differs from $V$ on a non-null set, it either changes
unanimously correct advice, contradicts itself, or gives a definite
recommendation on mixed problems. The first two possibilities have
positive penalty. In the third, the recommendation is bad at at least
one world at each changed problem. Since there are finitely many worlds,
some fixed world has positive integrated penalty. Thus $V$ weakly
improves every loss and strictly improves at least one: it dominates
every other model modulo null sets.

Define $C_p=V$ for every $p\in\Delta$. This selection is coherent and
undominated, and $C_p\preceq_iX$ for every $i,p,X$. In particular it is
competitive on every support. Each section
$\{p:C_p\preceq_iX\}$ is the whole simplex and is closed, proving (S).
All models are evaluable, proving (E); $V$ proves (LA); and a model tied
with $V$ at all worlds must equal it modulo null sets, proving (C).
The penalty table proves weak (G) and its stated strict parts, and
additivity proves (IB).
\end{proof}

This ranking system covers $V$ only. It does not establish full (A).
But if a class already covers other targets, adjoining this system
supplies coverage of $V$ without changing the coverage previously
obtained. The enlarged class satisfies the weakened guidance
requirement, not the original strict (G).

Two consequences of this change should be explicit. First, $D_i$ and
$V$ tie at world $i$, so the ideal model is no longer unique. The old
face lemma, which forces $C_{e_i}=D_i$ under strict (G), is unavailable;
the constant selection above illustrates its failure. Second, under
(IB), an indifference between good advice and silence transfers to
every common context. It cannot be made an exception only when one of
the two whole models happens to be $V$, while retaining strictness for
the identical local change elsewhere.

\section{Other boundary targets}

It would be premature to treat every boundary credal set alike.
Precise vertex models have zero lower probabilities at all but one
world, yet they are ordinary ideal models under strict (G). Zero lower
probability therefore does not, by itself, justify neutralizing the
value of good advice. Likewise, silence on some mixed problems is
common to ordinary imprecise models and does not by itself imply such
neutrality.

A useful distinction is between precise targets, imprecise targets
strictly inside the simplex, targets meeting its boundary, and targets
containing several vertices (including vacuous models on faces).
The many-world note's clique theorem concerns the last class. It does
not determine the coverage of the first three classes, or prove that
$V$ is the only excluded target. Removing $V$ from the coverage demand
removes this particular obstruction; a positive theorem for the
remaining class would still be needed.

\section{What survives of transfinite construction}

The original idea was to add the comparisons required by competitiveness
one at a time while preserving guidance improvement and the target's
undominatedness. There are two distinct issues: consistency when a
comparison is added, and preservation of closed sections at a limit.

\begin{lemma}\label{lem:union}
Fix a finite set of worlds, a common model domain, and an increasing
ordinal-indexed family of preorders $R_i^\alpha$ at each world. Suppose
each stage satisfies (G) and (IB), and a fixed target $M$ is undominated
at every stage. Then $R_i=\bigcup_\alpha R_i^\alpha$ is a preorder
satisfying (G) and (IB), and $M$ remains undominated. Any competitiveness
witness installed at a stage remains installed in the union.
\end{lemma}
\begin{proof}
Reflexivity is immediate. A finite transitivity chain lies in some
single stage, so transitivity persists. Every positive background
comparison persists. A reverse comparison forbidden by strict (G)
occurs at no stage and therefore not in the union. Context replacement
of any comparison occurs at the stage where that comparison occurs;
the inverse replacement gives the other direction of (IB).

If $X$ dominated $M$ in the union, the finitely many comparisons
$X R_i M$ would all hold at one stage. At a world where the union
comparison is strict, the reverse $M R_i X$ occurs at no stage.
Consequently $X$ would already dominate $M$ at that stage, a
contradiction. Competitiveness witnesses are positive comparisons,
which persist under extension.
\end{proof}

So the purely ordinal limit step has a useful positive answer.
\emph{Closedness does not follow.} Even with the selection fixed, an
increasing union of closed sections need not be closed:
\[
 \bigcup_{m\ge2}[0,1-1/m]=[0,1).
\]
Closing the union would add comparisons not justified by any stage.
After transitivity and (IB) are imposed, these may reverse a protected
strict guidance comparison or dominate the target. The missing
assertion is that some compatible way of closing them always exists;
the topological example alone does not prove that assertion false for
every possible construction.

\begin{proposition}[A sufficient invariant for closed sections]
Fix a selection $p\mapsto C_p$. Along an increasing construction,
suppose that whenever a comparison $C_p\preceq_iX$ is permanently
rejected, a relatively open neighbourhood $U$ of $p$ is reserved on
which every comparison $C_q\preceq_iX$ is permanently forbidden.
Suppose every such comparison is eventually either installed or
rejected, and no forbidden comparison is installed, including by
transitivity or (IB). Then the final comparison sections are closed.
\end{proposition}
\begin{proof}
If $p$ is outside the final section for $i,X$, its comparison was
rejected. The reserved neighbourhood $U$ remains outside that section.
Thus the section's complement is open.
\end{proof}

This replaces a proposed numerical ``uniform bound on lifts'' with a
topological invariant: permanent open regions of forbidden comparisons.
It does not yet supply a construction. One must prove an extension
lemma saying that the next competitiveness requirement, or the next
decision about a section, can be met while respecting all the previous
reservations after taking (IB) and transitive closure. Reserving a
neighbourhood at one world also requires competitiveness to be supplied
by other worlds throughout that neighbourhood. This is a genuine
compatibility obligation.

The earlier \emph{Ordinal Refinement by Slack Envelopes} proved that an
unrestricted refinement lemma is false, and instead established a
simultaneous refinement using extra structure. That success preceded
(IB); v3 subsequently showed why the benchmark/default-comparison
construction fails (IB). Neither result proves that all transfinite
constructions of incomplete rankings must fail. The new complete-case
theorem does exclude any construction that finishes with complete
rankings while retaining $V$ and the original axioms.

\section{Research priorities}

There are now three separate questions.
\begin{enumerate}[leftmargin=*]
\item With strict (G), (IB), (S), and complete rankings, characterize
which proper credal targets can be covered. Full coverage is impossible,
but its failure at $V$ does not answer this positive question.
\item If neutral abstention is appropriate for extreme targets, allow
the corresponding weak guidance ranking explicitly. Proposition
\ref{prop:neutral} handles $V$ completely under that revised requirement.
\item If full (A) and strict (G) are both retained, construct genuinely
incomplete rankings, or exclude them. A transfinite approach remains
possible in principle, but requires the compatibility extension lemma
just described. It must preserve the incomparabilities that the new
incomplete-ranking note identifies, rather than silently completing them.
\end{enumerate}
These questions should not be conflated. The first asks how broadly the
ordinary strict-guidance programme succeeds. The second changes how
abstention is evaluated. The third is the remaining existence problem
under the original full list of requirements.
\end{document}
